\section{Stress and strain equations for the hollow cylinder apparatus}
\label{sec:appendix_hca}

Table~\ref{tab:hca_equations} summarizes the stress and strain evaluation expressions used in the DEM-based hollow cylinder apparatus model. In the following, $R$ and $r$ denote the current outer and inner radii, $H$ and $H_0$ are the current and initial specimen heights, $T$ is the torque, $\theta$ is the twist angle, $p_i^{\prime}$, $p_o^{\prime}$, and $p_z^{\prime}$ are the radial and vertical effective pressures measured from particle--wall contact forces on the inner cylinder, outer cylinder, and top/bottom plates, respectively, $u_i$ and $u_o$ are the inner and outer radial displacements, and $f_{pw}^{i}$, $f_{pw}^{o}$, $f_{pw}^{t}$, $f_{pw}^{b}$ are particle--wall contact forces on the respective boundaries. In the constant-volume DEM approximation, all stresses computed from contact forces are effective stresses. The values of $\sigma_\theta^{\prime}$ and $\sigma_r^{\prime}$ are derived from equilibrium relations, $\varepsilon_z$ and $\gamma_{z\theta}$ are based on strain compatibility, and the remaining stress and strain expressions follow the assumption of linear elasticity \citep{Hight1983}.

\begin{table*}[t]
  \centering
  \caption{Equations of stress and strain in the hollow cylinder apparatus.}
  \label{tab:hca_equations}
  \begin{tabular}{@{}lll@{}}
    \toprule
    Component & Stress & Strain \\
    \midrule
    Inner pressure &
    $\displaystyle p_i^{\prime} = \frac{\sum\!\sqrt{{f_{pw_x}^{i}}^2+{f_{pw_y}^{i}}^2}}{2\pi r H}$ &
    --- \\[12pt]
    Outer pressure &
    $\displaystyle p_o^{\prime} = \frac{\sum\!\sqrt{{f_{pw_x}^{o}}^2+{f_{pw_y}^{o}}^2}}{2\pi R H}$ &
    --- \\[12pt]
    Vertical pressure &
    $\displaystyle p_z^{\prime} = \frac{\sum\!|f_{pw_z}^{t}|+\sum\!|f_{pw_z}^{b}|}{2\pi(R^2-r^2)}$ &
    --- \\[12pt]
    Vertical &
    $\sigma_z^{\prime} = p_z^{\prime}$ &
    $\displaystyle\varepsilon_z = -\frac{H-H_0}{H_0}$ \\[12pt]
    Circumferential &
    $\displaystyle\sigma_\theta^{\prime} = \frac{p_o^{\prime} R - p_i^{\prime} r}{R-r}$ &
    $\displaystyle\varepsilon_\theta = -\frac{u_o+u_i}{R+r}$ \\[12pt]
    Radial &
    $\displaystyle\sigma_r^{\prime} = \frac{p_o^{\prime} R + p_i^{\prime} r}{R+r}$ &
    $\displaystyle\varepsilon_r = -\frac{u_o-u_i}{R-r}$ \\[12pt]
    Shear &
    $\displaystyle\tau_{z\theta} = \frac{3T}{2\pi(R^3-r^3)}$ &
    $\displaystyle\gamma_{z\theta} = \frac{2\theta(R^3-r^3)}{3H(R^2-r^2)}$ \\[12pt]
    Major principal &
    $\displaystyle\sigma_1^{\prime} = \frac{\sigma_z^{\prime}+\sigma_\theta^{\prime}}{2}+\sqrt{\left(\frac{\sigma_z^{\prime}-\sigma_\theta^{\prime}}{2}\right)^{\!2}+\tau_{z\theta}^2}$ &
    $\displaystyle\varepsilon_1 = \frac{\varepsilon_z+\varepsilon_\theta}{2}+\sqrt{\left(\frac{\varepsilon_z-\varepsilon_\theta}{2}\right)^{\!2}+\left(\frac{\gamma_{z\theta}}{2}\right)^{\!2}}$ \\[12pt]
    Intermediate &
    $\sigma_2^{\prime} = \sigma_r^{\prime}$ &
    $\varepsilon_2 = \varepsilon_r$ \\[12pt]
    Minor principal &
    $\displaystyle\sigma_3^{\prime} = \frac{\sigma_z^{\prime}+\sigma_\theta^{\prime}}{2}-\sqrt{\left(\frac{\sigma_z^{\prime}-\sigma_\theta^{\prime}}{2}\right)^{\!2}+\tau_{z\theta}^2}$ &
    $\displaystyle\varepsilon_3 = \frac{\varepsilon_z+\varepsilon_\theta}{2}-\sqrt{\left(\frac{\varepsilon_z-\varepsilon_\theta}{2}\right)^{\!2}+\left(\frac{\gamma_{z\theta}}{2}\right)^{\!2}}$ \\[12pt]
    Deviatoric &
    $\displaystyle q = \sqrt{\tfrac{1}{2}\!\left[(\sigma_1^{\prime}-\sigma_2^{\prime})^2+(\sigma_2^{\prime}-\sigma_3^{\prime})^2+(\sigma_3^{\prime}-\sigma_1^{\prime})^2\right]}$ &
    $\displaystyle\varepsilon_q = \sqrt{\tfrac{2}{9}\!\left[(\varepsilon_1-\varepsilon_2)^2+(\varepsilon_2-\varepsilon_3)^2+(\varepsilon_3-\varepsilon_1)^2\right]}$ \\[6pt]
    \bottomrule
  \end{tabular}
\end{table*}
